% Source: https://artofproblemsolving.com/community/c3h1356843_dihedral_angle
% Stable identity: aops-topic-1356843-q01
Let the regular tetrahedron $ABCD$ have edge $2$ and let $M$ be the midpoint of edge $AB$. In the equilateral faces $ABC$ and $ABD$ the medians $CM$ and $DM$ are perpendicular to $AB$, so the angle $\theta=\angle CMD$ is the dihedral angle along $AB$ (it lies in the plane $CMD$, which is perpendicular to $AB$). Each of $CM$, $DM$ is an altitude of an equilateral triangle of side $2$, so $CM=DM=\sqrt3$, and $CD=2$.

\begin{asy}
usepackage("amsmath");
size(8cm);
defaultpen(fontsize(10pt));
pair pr(triple p){return (p.x+0.55*p.y, p.z+0.35*p.y);}
triple A3=(0,0,0),B3=(2,0,0),C3=(1,sqrt(3),0),D3=(1,sqrt(3)/3,sqrt(8/3)),M3=(1,0,0);
pair A=pr(A3),B=pr(B3),C=pr(C3),D=pr(D3),M=pr(M3);
draw(A--B--C--D--cycle,linewidth(0.9)); draw(B--D,linewidth(0.9)); draw(A--C,gray+dashed);
draw(M--D,red); draw(M--C,red+dashed);
dot(M);
label("$A$",A,SW); label("$B$",B,S); label("$C$",C,E); label("$D$",D,N); label("$M$",M,S);
pair o=(3.4,0);
pair Mp=o,Cp=o+(sqrt(3),0);
real ang=acos(1/3);
pair Dp=o+sqrt(3)*dir(degrees(ang));
draw(Mp--Cp--Dp--cycle,red);
draw(arc(Mp,0.35,0,degrees(ang)),red);
label("$\theta$",Mp+0.55*dir(degrees(ang)/2),red);
label("$M$",Mp,SW); label("$C$",Cp,SE); label("$D$",Dp,N);
label("$\sqrt{3}$",(Mp+Cp)/2,S); label("$\sqrt{3}$",(Mp+Dp)/2,NW); label("$2$",(Cp+Dp)/2,NE);
label("\mbox{}",(min(currentpicture,user=true)+max(currentpicture,user=true))/2,black); // keeps the text after the figure black when LaTeX typesets the labels inline
\end{asy}

By the law of cosines in triangle $CMD$ (drawn in true shape on the right),
\[
\cos\theta=\dfrac{CM^2+DM^2-CD^2}{2\cdot CM\cdot DM}=\dfrac{3+3-4}{2\cdot3}=\boxed{\dfrac{1}{3}},
\]
so $\theta=\arccos\dfrac{1}{3}\approx\boxed{70.5^\circ}$. All six edges are alike, so this is the dihedral angle along every edge.
